% AI as Computational Mathematician Rather Than Theorem Prover
% Claes Johnson. DRAFT.
\documentclass[11pt]{article}
\usepackage[margin=1.15in]{geometry}
\usepackage{amsmath,amssymb}
\usepackage{url}
\usepackage[hidelinks]{hyperref}
\setlength{\parskip}{3pt}

\newcommand{\grad}{\nabla}

\title{\vspace{-2em}AI as Computational Mathematician\\Rather Than Theorem Prover}
\author{Claes Johnson\\[0.3em]
\small KTH Royal Institute of Technology, Stockholm\\
\small \texttt{claesjohnson@gmail.com}}
\date{\today}

\begin{document}
\maketitle

\begin{abstract}
\noindent
AI has lately been demonstrated as a theorem prover: it produces analytical arguments and in some
cases formalises them in Lean. Its mathematical output includes substantial work in partial
differential equations, and almost no conventional numerical PDE --- no finite elements, no adaptive
computation, no error control, no computational turbulence. This paper asks whether that is the right
division of labour.

Numerical PDE already has its computational machinery, and there is no reason for a language model to
replace it. The comparative advantage lies one level up: \emph{AI decides what should be computed,
and conventional algorithms compute it}. Turbulent flow is the natural test case, because the
question that matters there is not whether the detailed solution is unique --- it is not --- but
whether mean-value outputs such as drag and lift are reproducibly computable. That question is open,
it is quantitative, and it is more naturally settled by controlled computation than by analytical
estimate. We argue that the role available to AI is that of a computational mathematician: designing
and running the experiments, interpreting them, and turning observed output uniqueness into
quantified and ultimately certified output uniqueness.
\end{abstract}

\section{A gap in the repository}\label{sec:gap}

OpenAI has recently demonstrated AI as a mathematical theorem prover, generating analytical arguments
and in some cases formalising or verifying them in Lean~\cite{openaiNS,openaiEuler}. The
mathematics produced includes hundreds of proof families and substantial work in partial differential
equations.

What is conspicuously thin is conventional numerical PDE: finite element and Galerkin methods, time
stepping, adaptive computation, a posteriori error control, computational turbulence. The machinery
by which partial differential equations are actually solved, and by which essentially every
engineering prediction involving them is made, is close to absent.

That could mean the subject is unsuitable for AI. We think it means the opposite, and that the
division of labour has been drawn in the wrong place.

\section{Numerical PDE already has its machinery}\label{sec:machinery}

For a system such as the incompressible Navier--Stokes equations, deterministic numerical algorithms
already perform the computation efficiently: finite element or Galerkin discretisation, time
stepping, nonlinear iteration, sparse linear algebra, multigrid and preconditioning, adaptive mesh
refinement. These are mature, they are fast, and their implementation is automated --- the FEniCS
project generates compiled code from a variational form, including the adjoint problem required for
error control~\cite{fenics}.

There is little reason for a language model to replace any of it. Once a numerical problem has been
specified, conventional algorithms and high-performance computing should perform the trillions of
floating-point operations; a model which attempted to do that arithmetic itself would be slower by
many orders of magnitude and less reliable.

The comparative advantage lies above the numerical algorithm:

\begin{center}
\emph{AI decides what should be computed; conventional algorithms compute it.}
\end{center}

\noindent The research loop is then
\[
\text{AI hypothesis}\ \rightarrow\ \text{PDE computation}\ \rightarrow\ \text{error and stability
analysis}\ \rightarrow\ \text{AI interpretation}\ \rightarrow\ \text{new computation},
\]
in which AI writes and modifies code, chooses numerical experiments, varies parameters, launches
ensembles, identifies unexpected behaviour, decides where refinement is required, compares methods,
analyses results and formulates the next experiment. It might also discover improved numerical
methods --- stabilisations, preconditioners, adaptive criteria, reduced models, error estimators.
Once discovered and validated, these would become ordinary deterministic algorithms, which is the
proper fate of a good method.

\section{Turbulence as the test case}\label{sec:turb}

Computational turbulence is where this division of labour can be put to work, and the reason is that
the interesting question there is not the one analysis has been asking.

\paragraph{Three solution concepts.} The history is instructive. A \emph{pointwise} solution of the
Euler or Navier--Stokes equations exists and is unstable, so determines nothing, since it does not
occur. Leray's \emph{weak} solution of 1934~\cite{leray} exists globally and is not unique, so
determines nothing either, since it does not say which. The concept developed in~\cite{book} is a
weak solution together with \emph{control of the residual}, and what it delivers is different in
kind: the flow field is not determined pointwise in space and time, and mean-value outputs are.

\paragraph{Output uniqueness.} Suppose $U_h$ and $V_h$ are approximate turbulent solutions computed
with different meshes, perturbations or schemes. It is entirely possible that
\[
\|U_h-V_h\| \not\rightarrow 0 ,
\qquad\text{while for a mean-value output }M,\qquad
|M(U_h)-M(V_h)| \rightarrow 0 .
\]
This is \emph{output} or \emph{weak} uniqueness: the detailed turbulent solution need not be uniquely
predictable for physically relevant mean quantities to be uniquely predictable. Stabilised Galerkin
methods together with primal/dual a posteriori error analysis were developed precisely to investigate
it, and the computations indicate stability of mean outputs even where the underlying flow is
unstable~\cite{book,hjv,hjja}.

That is the phenomenon. Whether it holds in general, and under what conditions, is open.

\section{Why an analytical proof need not be the goal}\label{sec:proof}

A traditional programme would now seek a theorem establishing this stability. It is worth asking why
an analytical proof should be considered superior, because three things argue against it here.

\paragraph{Worst-case constants.} Analytical estimates in turbulent flow are worst-case by
construction, and the constants they produce grow extravagantly with the Reynolds number. A bound
which is true and of size $\exp\exp\exp A$ is not a bound any flow can be said to respect, and no
measurement could distinguish it from no bound at all. Computation of the actual dual dynamics may
reveal stability constants smaller by many orders of magnitude --- and those are the numbers that
decide whether an output is determined in practice.

\paragraph{No structure to exploit.} Symbolic analysis proceeds by rewriting, and rewriting needs
something to grip: a symmetry, a similarity variable, an ansatz. Every construction of a singular
solution buys its tractability with one. Turbulence affords none --- it is defined by carrying
structure at every scale with no organising principle to exploit --- and that is not a gap in
anyone's technique but what the word means.

\paragraph{What is actually unresolved.} The essential issue is not the absence of a traditional
proof. It is \emph{convincing control of the limiting process $h\rightarrow0$}. That control could
itself be computational, and this is the point on which the paper turns.

\section{Massive computational experiment}\label{sec:experiment}

Here AI could do what a human research group cannot, which is to run the investigation at scale.

Thousands of turbulent computations, varying independently the mesh size $h$, the time step $\Delta
t$, the viscosity $\nu$, the mesh geometry, the stabilisation, the numerical method, the initial
perturbations and the boundary approximation. A systematic search through those computations for
observables satisfying
\[
M(U_h)\ \longrightarrow\ M^*
\]
even where the instantaneous solutions differ substantially. Primal and dual problems computed
together, so that residuals and a posteriori estimates quantify the uncertainty in each output
rather than leaving it to be asserted. Independent methods applied to the same case, to test whether
the same output is selected by machinery that shares no code.

This is laborious, it is not deep, and it is exactly the kind of work an AI system is suited to:
hypothesis, experiment, interpretation, next experiment, at a rate and breadth no group sustains.

The progression available is then from \emph{computational evidence}, to \emph{quantified
computational evidence}, and possibly to \emph{certified computational mathematics} --- rigorous
residual estimates, interval arithmetic, verified numerics, turning parts of the computation into
computer-assisted proof. Formal theorem proving need not be the starting point. It may be the
destination.

\section{Why certification is reachable here}\label{sec:certify}

There is a structural reason for optimism about that last step, and it distinguishes the computational
route from the symbolic one.

Neither a long computation nor a long proof can be verified step by step by a person: one performs
billions of floating-point operations, the other contains more inferences than anyone will read. But
they escape that position differently, and only one of them escapes.

A numerical result can be checked \emph{without checking the computation}. Substitute the computed
field into the equations and evaluate the residual; weight it by the solution of an adjoint problem
and it bounds the error in any output of interest. The decisive feature is that this test is
independent of how the field was obtained --- it inspects no mesh, no solver, no arithmetic, and
would be the same test applied to a candidate that arrived by guess. Nobody therefore needs to verify
the individual steps, and nobody does: an error at the three-millionth operation is caught because
the residual comes out large. The equation audits the answer, once, wholesale.

A proof admits no such operation. What it produces is a statement rather than an object, and a
statement cannot be substituted into anything; its correctness simply \emph{is} the correctness of
its steps, which must be checked one at a time. A proof assistant does this mechanically and
reliably --- and establishes only that the steps follow from the stated definitions, not that the
definitions say what was intended, which cannot be mechanised.

So the certification of computational results is not an aspiration borrowed from the symbolic side.
It is the native mode of the subject, and the a posteriori error estimate is already a theorem about
the exact solution, obtained by computation.

\section{Two paradigms}\label{sec:paradigms}

The contrast can be put compactly. The first is currently prominent:
\[
\boxed{\ \text{AI}\ \rightarrow\ \text{analytical proof}\ \rightarrow\ \text{formal verification}\ }
\]
The second is the one proposed here:
\[
\boxed{\ \text{AI}\ \rightarrow\ \text{numerical PDE}\ \rightarrow\ \text{computed solutions}\
\rightarrow\ \text{error control}\ \rightarrow\ \text{scientific conclusion}\ }
\]

In the second, AI functions less as a theorem prover than as a \emph{computational mathematician} or
experimental scientist. The numerical solver is the laboratory; AI designs the experiments, controls
them, interprets the results, searches for reproducible mathematical structure, and decides what to
do next.

Two observations about the comparison, neither of them rhetorical.

\paragraph{The second paradigm already carries its error control; the first does not.} A finite
element computation returns an output together with a bound on how wrong that output may be, computed
from the residual weighted by an adjoint solution. A language model returns an argument with no
bound of any kind, and the problem it now faces --- an output too long to check and too important to
trust blindly --- is precisely the one numerical analysis met when machines first outran hand
calculation, and answered by computing the error alongside the answer. That answer transfers in
principle, and the machinery is numerical in both cases, since a model is itself a numerical object.

\paragraph{And the synthesis is not new.} Developing analysis and computation together, rather than
treating computation as a service added afterwards, has been pursued as a programme for some
twenty-five years under the name \emph{Applied Mathematics: Body and Soul}~\cite{bodysoul,bsbooks},
where the name is the thesis: the body of a mathematical subject is its computational content and the
soul its analytical content, and a subject which has lost either is crippled. What is new is a
machine capable of operating the loop.

\section{The question for turbulence}\label{sec:question}

For turbulent Navier--Stokes flow the scientifically important question may not be whether every
detailed weak solution is unique. It is not, and the convex-integration constructions have settled
that in the negative. The question is whether physically relevant mean quantities are reproducibly
and accurately computable --- which is what an aircraft designer asks, what a wind tunnel measures,
and what the benchmark workshops test by requiring predictions before the measurements are
released~\cite{cs1,ifasd}.

So the challenge for AI becomes:

\begin{center}
\emph{Can AI turn observed computational output uniqueness into quantitatively controlled --- and
ultimately certified --- computational output uniqueness?}
\end{center}

\noindent That is a different contribution from producing more analytical proofs, and on the evidence
of the subject's history it is the more valuable one. The equations have been available since 1755
and the analytical programme has not delivered a drag coefficient. The computational programme has,
and what it lacks is not power but the systematic exploration of its own limiting behaviour --- which
is exactly the work now within reach.

\begin{thebibliography}{99}

\bibitem{openaiNS} OpenAI, \emph{Finite time blowup for Navier--Stokes}, preprint, 2026.
\newblock \url{https://cdn.openai.com/pdf/32d9f210-8b73-45e0-91bc-82a30aef8a9a/navier-stokes.pdf}

\bibitem{openaiEuler} OpenAI, \emph{Finite time blowup for the Euler equation}, preprint,
8~September 2026, 57~pp.
\newblock \url{https://cdn.openai.com/pdf/315b36cd-ec98-4023-8342-93345194ece1/euler.pdf}

\bibitem{leray} J.~Leray, \emph{Sur le mouvement d'un liquide visqueux emplissant l'espace},
Acta Math.\ \textbf{63} (1934) 193--248.

\bibitem{book} J.~Hoffman and C.~Johnson, \emph{Computational Turbulent Incompressible Flow},
Applied Mathematics: Body and Soul, vol.~4, Springer, 2007.

\bibitem{hjv} J.~Hoffman, J.~Jansson and R.~Vilela de Abreu, \emph{Adaptive modeling of turbulent
flow with residual based turbulent kinetic energy dissipation}, Comput.\ Methods Appl.\ Mech.\
Engrg.\ \textbf{200} (2011) 2758--2767.

\bibitem{hjja} J.~Hoffman, J.~Jansson, N.~Jansson and R.~Vilela de Abreu, \emph{Towards a
parameter-free method for high Reynolds number turbulent flow simulation based on adaptive finite
element approximation}, Comput.\ Methods Appl.\ Mech.\ Engrg.\ \textbf{288} (2015) 60--74.

\bibitem{fenics} A.~Logg, K.-A.~Mardal and G.~N.~Wells (eds.), \emph{Automated Solution of
Differential Equations by the Finite Element Method}, Lecture Notes in Computational Science and
Engineering~84, Springer, 2012.

\bibitem{bodysoul} The Body and Soul project, \url{www.bodysoulmath.org}.

\bibitem{bsbooks} K.~Eriksson, D.~Estep and C.~Johnson, \emph{Applied Mathematics: Body and Soul},
vols.~1--3, Springer, 2004.

\bibitem{cs1} J.~Jansson, \emph{Adaptive Direct FEM simulation with Unicorn/FEniCS-HPC for CS1},
5th International Workshop on High Order CFD Methods; Zenodo, 2022.
\newblock \url{https://doi.org/10.5281/zenodo.6942599}

\bibitem{ifasd} J.~Jansson, \emph{Adaptive Euler: validation against NASA High Lift Prediction
Workshop~4, and preliminary tail-buffet and transonic-buffet simulations}, presented at IFASD~2026.

\end{thebibliography}

\end{document}
